\begin{ex}
Số gồm 8 chục nghìn, 5 nghìn, 3 trăm, 4 chục và 2 đơn vị viết là gì?
\shortans{85342}
\loigiai{
Số đó là 85342.
}
\end{ex}

\begin{ex}
Số 74 520 có chữ số hàng chục nghìn là bao nhiêu?
\shortans{7}
\loigiai{
Chữ số hàng chục nghìn là 7.
}
\end{ex}

\begin{ex}
Số liền sau của số 99 999 là số nào?
\shortans{100000}
\loigiai{
$99999 + 1 = 100000$.
}
\end{ex}

\begin{ex}
Số liền trước của số 70 000 là số nào?
\shortans{69999}
\loigiai{
$70000 - 1 = 69999$.
}
\end{ex}

\begin{ex}
Làm tròn số 48 650 đến hàng nghìn ta được số nào?
\shortans{49000}
\loigiai{
Làm tròn lên 49 000.
}
\end{ex}

\begin{ex}
Làm tròn số 73 200 đến hàng chục nghìn ta được số nào?
\shortans{70000}
\loigiai{
Làm tròn thành 70 000.
}
\end{ex}

\begin{ex}
Số La Mã XIX biểu thị số tự nhiên nào?
\shortans{19}
\loigiai{
XIX là 19.
}
\end{ex}

\begin{ex}
Số La Mã XIV biểu thị số tự nhiên nào?
\shortans{14}
\loigiai{
XIV là 14.
}
\end{ex}

\begin{ex}
Trong các số $45600, 46500, 54600, 56400$, số nào bé nhất?
\shortans{45600}
\loigiai{
45600 là số bé nhất.
}
\end{ex}

\begin{ex}
Trong các số $82450, 84250, 82540, 84520$, số nào lớn nhất?
\shortans{84520}
\loigiai{
84520 là số lớn nhất.
}
\end{ex}

\begin{ex}
Số lớn nhất có 5 chữ số khác nhau là số nào?
\shortans{98765}
\loigiai{
Số đó là 98765.
}
\end{ex}

\begin{ex}
Số bé nhất có 5 chữ số khác nhau là số nào?
\shortans{10234}
\loigiai{
Số đó là 10234.
}
\end{ex}

\begin{ex}
Số 55 555 có bao nhiêu chữ số 5?
\shortans{5}
\loigiai{
Có 5 chữ số 5.
}
\end{ex}

\begin{ex}
Viết số thích hợp: $50000 + 7000 + 200 + 4 = \dots$
\shortans{57204}
\loigiai{
Số đó là 57204.
}
\end{ex}

\begin{ex}
Khoảng cách giữa hai số 60 000 và 58 500 là bao nhiêu?
\shortans{1500}
\loigiai{
$60000 - 58500 = 1500$.
}
\end{ex}

\begin{ex}
Có bao nhiêu số tròn chục nghìn từ 10 000 đến 100 000?
\shortans{10}
\loigiai{
10 số.
}
\end{ex}

\begin{ex}
Số La Mã XX tương ứng số tự nhiên nào?
\shortans{20}
\loigiai{
XX là 20.
}
\end{ex}

\begin{ex}
Đồng hồ La Mã chỉ vào số XII nghĩa là mấy giờ?
\shortans{12}
\loigiai{
XII là 12 giờ.
}
\end{ex}

\begin{ex}
Làm tròn số 850 đến hàng trăm được số nào?
\shortans{900}
\loigiai{
Làm tròn lên 900.
}
\end{ex}

\begin{ex}
Làm tròn số 342 đến hàng chục được số nào?
\shortans{340}
\loigiai{
Làm tròn thành 340.
}
\end{ex}

\begin{ex}
Điền số thích hợp: $79999 < \dots < 80001$.
\shortans{80000}
\loigiai{
Số ở giữa là 80000.
}
\end{ex}

\begin{ex}
Giá trị của chữ số 6 trong số 65 432 là bao nhiêu?
\shortans{60000}
\loigiai{
Chữ số 6 ở hàng chục nghìn có giá trị 60 000.
}
\end{ex}

\begin{ex}
Hiệu của số lớn nhất có 5 chữ số và số bé nhất có 5 chữ số là bao nhiêu?
\shortans{89999}
\loigiai{
$99999 - 10000 = 89999$.
}
\end{ex}

\begin{ex}
Số 100 000 đọc là một trăm nghìn, có bao nhiêu chữ số 0?
\shortans{5}
\loigiai{
Có 5 chữ số 0.
}
\end{ex}

\begin{ex}
Tìm $x$ là số tròn nghìn biết $68000 < x < 70000$.
\shortans{69000}
\loigiai{
Số đó là 69000.
}
\end{ex}
